\documentclass[11pt]{ctexart}
\usepackage[a4paper,margin=1in]{geometry}
\usepackage{enumitem}
\usepackage{hyperref}
\usepackage{xcolor}
\usepackage{titlesec}

\hypersetup{
    colorlinks=true,
    linkcolor=blue,
    urlcolor=blue
}

\setlist[itemize]{leftmargin=1.35em,itemsep=0.12em,topsep=0.2em}
\setlist[enumerate]{leftmargin=1.55em,itemsep=0.12em,topsep=0.2em}
\setlength{\parskip}{0.34em}
\linespread{0.98}
\titlespacing*{\section}{0pt}{0.8em}{0.32em}
\titlespacing*{\subsection}{0pt}{0.55em}{0.18em}

\title{海报展示讲稿\\\large 通信受限时的分布式影响力最大化}
\author{陈亦新 \and 张弼弛 \and 叶瑶勤}
\date{CS240 算法设计与分析课程，2026 年春季\\2026 年 6 月 18 日}

\begin{document}
\maketitle

\section*{使用方式}

这份讲稿是给面对海报做展示时使用的，不是幻灯片讲稿。推荐讲的时候按海报从左到右移动：
\begin{enumerate}
    \item 先从顶部的研究问题开始；
    \item 再解释左侧的优化目标、扩散模型和子模性；
    \item 接着讲中间的算法路线和分布式流程；
    \item 最后用右侧的大规模证据、通信折中和主要结论收尾。
\end{enumerate}

下面的主讲稿大约 4--5 分钟，适合正式展示；短版讲稿大约 90 秒，适合老师或同学只是快速走到海报前询问时使用。后面的问答部分可以作为现场答辩准备。

\section*{正式展示讲稿：4--5 分钟}

\subsection*{1. 开场}

大家好，我这个项目研究的是影响力最大化问题，重点不是只看“能不能找到最好的种子节点”，而是看当通信成为瓶颈时，我们还能保留多少影响力质量。

在社交网络、信息传播和推荐系统里面，一个常见任务是先选择少量初始用户，然后希望信息通过网络扩散到尽可能多的节点。经典影响力最大化通常假设我们可以集中访问整张图，并反复评估候选节点。但在真实系统中，图可能分布在多个机器或多个社区里，集中式计算会带来通信和计算成本。所以我在这个项目里比较了集中式贪心、CELF 加速贪心，以及一种 GreeDI 风格的分布式候选共享方法。

\subsection*{2. 问题与扩散模型}

这里的优化目标是：给定图 $G=(V,E)$、预算 $k$，选择一个种子集合 $S$，使独立级联模型下的期望传播规模 $\sigma(S)$ 最大。

我采用的是独立级联模型，也就是 IC 模型。每个被激活的节点只有一次机会去激活它的邻居。边上的传播概率使用加权级联设定：如果边指向节点 $v$，那么传播概率是 $1/\mathrm{indegree}(v)$。这个设定的好处是它和网络结构相关，不需要额外人工指定每条边的概率。

从理论上看，IC 模型下的影响力函数是单调且子模的，所以经典贪心算法可以达到 $1-1/e$ 的近似保证。但是这个保证的代价是需要大量传播模拟。在大图上，这个成本很快会变成主要瓶颈。

\subsection*{3. 算法路线}

海报中间这部分展示的是算法路线。

第一类是简单基线，包括随机选择、加权出度和 PageRank。它们几乎不需要传播模拟，所以非常快，但不一定直接优化影响力函数。

第二类是集中式蒙特卡洛贪心。它每一轮都会枚举候选节点，估计加入该节点后的边际增益，然后选择增益最大的节点。它是最直接的理论基准，但运行时间最长。

第三类是 CELF，也叫懒惰贪心。CELF 利用子模性：节点的边际增益只会下降，所以很多候选节点不需要在每一轮都重新计算。它保留了和贪心非常接近的选择逻辑，但大幅减少传播估计次数。

最后是分布式候选共享方法。做法是把图划分给多个工作节点，每个工作节点先在本地选出少量候选种子，然后只把这些候选节点传回中心，中心再在候选池上做最终选择。这里的两个关键参数是工作节点数量 $p$ 和每个工作节点传回的候选数量 $qk$。这个方法的核心问题就是：通信变少以后，影响力质量会下降多少？

\subsection*{4. 实验设置}

实验图包括 Erdős--Rényi 随机图、Barabási--Albert 无标度图、Watts--Strogatz 小世界图，以及 Karate Club 真实小图。预算 $k$ 随图规模变化，小图和大图都有测试。

评估时，我把最终种子集合重新放回完整图上，用相同的 IC 模型做传播估计。这样可以避免一个方法只是在自己的局部视角上看起来很好。比较指标包括传播规模、相对于集中式贪心的比例、运行时间、传播估计次数，以及分布式方法需要上传的候选数量。

\subsection*{5. 主要结果}

第一个重要结果是，CELF 是最稳定的工程改进。在大规模预设中，它和贪心的平均传播质量几乎一样，平均比例大约是 $0.99$，但传播估计次数减少了约 $74.1\%$，运行时间减少了约 $83.5\%$。这说明如果目标是在不牺牲质量的前提下降低计算成本，CELF 是非常清楚的选择。

第二个结果是，分布式候选共享可以在较低通信量下保持不错的影响力质量。比如当 $p=2,q=1$ 时，平均传播规模和集中式贪心几乎相同，同时每个实例只需要上传大约 19 个候选节点。随着 $p$ 和 $q$ 改变，质量并不是严格单调上升的，因为图划分会打断跨分区传播结构，候选池变大也会带来更多冗余和模拟噪声。

第三个有意思的现象是，加权出度在一些平均结果上甚至超过了蒙特卡洛贪心。这不是和贪心理论矛盾，因为理论保证针对的是精确边际增益，而我的实验中贪心使用有限次蒙特卡洛模拟；同时加权级联模型本身也会让高度节点或指向低入度邻居的节点非常有优势。所以这个结果更多说明：在这种概率设定和有限模拟预算下，简单结构启发式可以很强。

\subsection*{6. 结论}

总结来说，这个项目的结论有三点。

第一，集中式贪心是理论上最自然的基准，但它的模拟成本很高。第二，CELF 基本保留了贪心质量，同时显著降低计算量。第三，分布式候选共享展示了影响力质量和通信成本之间的折中：只传少量本地候选也能得到接近集中式的结果，但参数选择和图结构会明显影响表现。

如果继续做下去，我会尝试更大的真实网络、更多划分策略，以及像 TIM、IMM 或 PaC-IM 这类更强的影响力最大化算法，把它们和这里的通信受限设置进行更公平的比较。

\section*{短版讲稿：约 90 秒}

我这个项目研究的是通信受限条件下的影响力最大化。经典问题是：在一个网络里选择 $k$ 个种子节点，让信息在独立级联模型下传播到尽可能多的节点。理论上，IC 模型下的影响力函数是单调子模的，所以贪心算法有 $1-1/e$ 的近似保证。但贪心需要大量蒙特卡洛传播模拟，在大图或分布式图上会很贵。

所以我比较了四类方法：随机、加权出度、PageRank 这样的快速基线；集中式蒙特卡洛贪心；利用子模性的 CELF 懒惰贪心；以及 GreeDI 风格的分布式候选共享方法。分布式方法让每个工作节点先在本地选候选，只上传少量节点，再由中心完成最终选择。

主要结果是：CELF 几乎保持了贪心的传播质量，但在大规模预设里减少了约 $74.1\%$ 的传播估计次数和约 $83.5\%$ 的运行时间。分布式候选共享在通信很少时也能接近集中式贪心，比如 $p=2,q=1$ 时平均质量接近贪心，只需要上传约 19 个候选节点。整体结论是，在影响力最大化中，通信和计算成本不能只当作实现细节；它们会直接改变我们应该选择的算法。

\section*{常见问题与回答}

\subsection*{问题 1. 为什么选择影响力最大化这个问题？}

因为它是算法设计里很经典的问题，同时又有很强的现实背景，比如社交网络传播、营销、推荐和信息扩散。它也很适合展示理论和系统约束之间的关系：理论上贪心有近似保证，但真实系统里计算和通信成本同样重要。

\subsection*{问题 2. 为什么使用独立级联模型？}

IC 模型简单、标准，而且在影响力最大化文献中使用非常广泛。它把传播过程建模成每条边一次独立尝试，所以便于用蒙特卡洛模拟估计期望传播规模。更重要的是，在 IC 模型下影响力函数具有单调性和子模性，这让贪心算法有理论依据。

\subsection*{问题 3. 加权级联概率 $1/\mathrm{indegree}(v)$ 有什么含义？}

它表示一个入度很大的节点同时受到很多邻居影响，所以单个邻居成功激活它的概率较小；入度小的节点则更容易被某个邻居激活。这样每条边的概率自动由图结构决定，不需要额外调参。

\subsection*{问题 4. 为什么不用固定传播概率，比如每条边都是 0.1？}

固定概率也可以用，但会弱化不同图结构之间的差异。加权级联模型更能反映目标节点入度对传播难度的影响，也更接近很多影响力最大化实验里的标准设置。

\subsection*{问题 5. 贪心算法既然有理论保证，为什么还需要 CELF？}

理论保证说明贪心的解质量好，但不说明它便宜。普通贪心每一轮都要重新评估大量候选节点，模拟次数很高。CELF 利用子模性缓存边际增益，避免很多不必要的重新计算，所以它主要解决的是效率问题。

\subsection*{问题 6. CELF 会改变贪心解吗？}

如果边际增益被精确计算，CELF 和普通贪心会选择相同的种子集合，因为它只是改变评估顺序，不改变贪心准则。在我的实验中边际增益来自蒙特卡洛估计，所以可能有轻微差异，但平均传播质量非常接近。

\subsection*{问题 7. 为什么加权出度有时能超过贪心？这是否说明贪心失效？}

不说明贪心失效。贪心的理论保证建立在精确边际增益上，而实验中使用的是有限次数蒙特卡洛模拟，会有估计噪声。另外，在加权级联模型下，结构启发式本身可能非常适合某些图，所以加权出度在平均值上超过模拟贪心并不矛盾。

\subsection*{问题 8. 分布式候选共享里的 $p$ 和 $q$ 分别是什么意思？}

$p$ 是工作节点数量，也就是图被分成多少个部分；$q$ 控制每个工作节点上传多少候选节点。具体来说，每个工作节点会上传大约 $qk$ 个本地候选。$p$ 越大，局部视角越碎；$q$ 越大，通信量越高，但中心可选的候选池也越大。

\subsection*{问题 9. 为什么分布式方法中 $q$ 增大后效果不一定更好？}

理论上候选池更大应该提供更多选择，但实验里有两个影响因素。第一，图划分会破坏跨分区传播结构，本地好的节点不一定全局最好。第二，蒙特卡洛估计有噪声，候选池变大后也可能引入更多冗余节点。因此质量不一定严格单调提升。

\subsection*{问题 10. 你如何保证不同算法比较公平？}

最终评估时，所有算法选出的种子集合都会放回同一张完整图上，用相同的 IC 模型和相同的评估流程重新估计传播规模。因此分布式方法不是只在自己的局部子图上评估，而是接受完整图上的统一评估。

\subsection*{问题 11. 这个项目的主要局限是什么？}

主要局限有三个。第一，图规模还比较小，尤其和真实社交网络相比。第二，使用的是合成图加 Karate Club 小图，真实大规模网络还不够。第三，实验重点是贪心、CELF 和 GreeDI 风格方法，还没有实现 TIM、IMM、GreediRIS 或 PaC-IM 等更先进算法。

\subsection*{问题 12. 如果继续做，你会怎么扩展？}

我会从三个方向扩展：第一，使用更大的真实网络数据；第二，比较更多图划分策略，比如随机划分、社区划分和边界感知划分；第三，引入更先进的反向可达集合方法，比如 TIM 或 IMM，再看它们在通信受限环境下是否仍然保持优势。

\section*{备用数据点}

\begin{itemize}
    \item CELF 在大规模预设中平均传播比例约为 $0.99$，非常接近集中式贪心。
    \item CELF 相比普通贪心减少约 $74.1\%$ 的传播估计次数。
    \item CELF 相比普通贪心减少约 $83.5\%$ 的实际运行时间。
    \item 分布式设置 $p=2,q=1$ 的平均质量接近集中式贪心，平均上传候选数量约为 19。
    \item 最核心的一句话总结：这个项目说明了在影响力最大化中，通信和计算成本会直接影响算法选择，而不仅仅是实现细节。
\end{itemize}

\end{document}
